\subsection{Dostupnost modelu}
\label{subsec:criteria-availability}

Dostupností modelu se rozumí soubor podmínek, které určují, zda a jakým způsobem lze model v systému Kelvin reálně provozovat.
Nejde pouze o technický přístup k modelu, ale i o právní a provozní aspekty, které mohou některé varianty zcela vyloučit:

\begin{enumerate}
    \item \textbf{Způsob přístupu} -- model může být dostupný jako cloudové API, kdy veškerý výpočet zajišťuje externí poskytovatel, nebo jako open-source model ke stažení a lokálnímu nasazení.
    Cloudové API nabízí snadnou integraci bez hardwarových nároků, ale zavádí závislost na externím poskytovateli.
    Open-source modely naopak umožňují lokální nasazení a plnou kontrolu nad zpracováním dat.

    \item \textbf{Licenční podmínky} -- pro akademické použití v rámci vysoké školy je nutné ověřit, zda licence modelu takové nasazení umožňuje.
    Některé modely jsou dostupné pod permisivními licencemi (Apache~2.0, MIT), jiné však obsahují omezení pro komerční použití nebo vyžadují uzavření samostatné licenční smlouvy.

    \item \textbf{Geografická dostupnost a datová centra} -- jak bylo rozebráno v \secref[sekci]{subsubsec:gdpr-data-transfer} o přenosu osobních údajů, pro dodržení GDPR je důležité, aby poskytovatel cloudového API provozoval datová centra v regionech, které jsou s~GDPR v souladu.
    Tato podmínka se týká výhradně cloudového nasazení, protože u on-premise řešení data školu neopouštějí.

    \item \textbf{Stabilita a podpora} -- u cloudového API je klíčová spolehlivost služby a předvídatelnost cenového modelu.
    Změna API, zrušení konkrétního modelu nebo úprava cenových podmínek ze strany poskytovatele mohou ohrozit funkčnost celého systému.
    Pro školní prostředí, kde je systém provozován dlouhodobě, je rovněž důležitá dostupnost dokumentace a garantovaná doba odezvy (SLA).
\end{enumerate}

\endinput